% !TeX program = xelatex
%% =====================================================================
%%  [เฉลย] ใบกิจกรรมในชั้นเรียน (Worksheet Solution) บทที่ 3 : ความสัมพันธ์และฟังก์ชัน --- สัปดาห์ที่ 9
%%  เอกสารเฉลยสำหรับผู้สอน: 2 หน้าเต็ม (ฉบับสมบูรณ์)
%% =====================================================================
\documentclass[a4paper,12pt]{article}
\usepackage{fontspec}
\usepackage{amsmath,amssymb}
\usepackage[margin=1.2cm,top=1.0cm,bottom=1.0cm]{geometry}
\usepackage{enumitem}
\usepackage{array}
\usepackage{xcolor}
\usepackage{tikz}

\XeTeXlinebreaklocale "th_TH"
\XeTeXlinebreakskip = 0pt plus 1pt
\setmainfont[Script=Thai,Scale=1.15]{TH Sarabun New}
\definecolor{navy}{HTML}{1F3A93}
\definecolor{mint}{HTML}{E8F5E9}
\definecolor{lightnavy}{HTML}{EAEEF7}
\definecolor{ans}{HTML}{C0392B} % สีเฉลยเน้นชัดเจน
\newcommand{\ans}[1]{\textbf{\color{ans}#1}}

\renewcommand{\arraystretch}{1.15}
\setlist{topsep=1pt,itemsep=2pt,parsep=0pt}
\pagestyle{empty}

\begin{document}

%% ---------- หัวกระดาษหน้า 1 ----------
\noindent
\begin{minipage}[c]{0.78\textwidth}
  {\Large\bfseries\color{navy} [เฉลย] ใบกิจกรรมในชั้นเรียน บทที่ 3 : ความสัมพันธ์และฟังก์ชัน (สัปดาห์ที่ 9)}\\[0.1em]
  {\normalsize รายวิชา 4091606 คณิตศาสตร์สำหรับคอมพิวเตอร์}\\[0.5em]
  {\bfseries\color{ans} ฉบับเฉลยและแนวทางคำตอบสำหรับผู้สอน (Master Answer Key)}
\end{minipage}\hfill
\begin{minipage}[c]{0.2\textwidth}
  \centering
  \fbox{\parbox{2.8cm}{\centering\bfseries\color{navy}เฉลยประจำสัปดาห์\\ Week 09 KEY}}
\end{minipage}
\par\vspace{0.3em}\hrule\vspace{0.4em}

\noindent\textbf{คำชี้แจง:} เอกสารเฉลยคำตอบพร้อมวิธีทำอย่างละเอียดทุกขั้นตอน สำหรับผู้สอนใช้อภิปรายและประเมินผลการเรียนรู้
\vspace{0.3em}

%% ---------- กิจกรรม 1 ----------
\noindent\textbf{กิจกรรมที่ 1 : ลำดับเลขคณิตและเรขาคณิตในชีวิตประจำวัน (Arithmetic \& Geometric Sequences)} \hfill \textit{(8 นาที)}
\begin{enumerate}[label=\arabic*)]
  \item \textbf{การออมเงินแบบฝากเพิ่มรายเดือน (ลำดับเลขคณิต):} เปิดบัญชีเริ่มต้น 1,000 บาท ($a_1 = 1,000$) และฝากเพิ่มเดือนละ 500 บาทเท่ากันทุกเดือน ($1000, 1500, 2000, \dots$)
        \begin{enumerate}[label=(\alph*),itemsep=1pt]
          \item ผลต่างร่วม ($d$) = \ans{500 บาท} (ได้จาก $1,500 - 1,000 = 500$)
          \item เขียนสูตรพจน์ทั่วไป $a_n = a_1 + (n-1)d = 1,000 + (n-1)(500) =$ \ans{$500n + 500$}
          \item เมื่อครบ 1 ปี (สิ้นเดือนที่ 12) จะมียอดเงินในบัญชี $a_{12} = 500(12) + 500 =$ \ans{6,500 บาท}
        \end{enumerate}
  \item กำหนดลำดับเลขคณิต $3, 7, 11, 15, \dots$
        \begin{enumerate}[label=(\alph*),itemsep=1pt]
          \item ผลต่างร่วม ($d$) = \ans{$7 - 3 = 4$}
          \item พจน์ทั่วไป ($a_n$) = $3 + (n-1)(4) =$ \ans{$4n - 1$}
          \item พจน์ที่ 50 ($a_{50}$) = $4(50) - 1 =$ \ans{$199$}
        \end{enumerate}
  \item \textbf{การเติบโตของเงินลงทุนแบบทบเท่าตัว (ลำดับเรขาคณิต):} เงินต้น 10,000 บาท ($a_1 = 10,000$) มูลค่าเพิ่มขึ้นเป็น 2 เท่าทุกปี ($10000, 20000, 40000, \dots$)
        \begin{enumerate}[label=(\alph*),itemsep=1pt]
          \item อัตราส่วนร่วม ($r$) = \ans{2} (ได้จาก $\frac{20,000}{10,000} = 2$)
          \item เขียนสูตรพจน์ทั่วไป $a_n = a_1 \cdot r^{n-1} =$ \ans{$10,000 \cdot 2^{n-1}$}
          \item สิ้นปีที่ 5 เงินลงทุนจะมีมูลค่า $a_5 = 10,000 \cdot 2^4 = 10,000 \times 16 =$ \ans{160,000 บาท}
        \end{enumerate}
  \item กำหนดลำดับเรขาคณิต $2, 6, 18, 54, \dots$
        \begin{enumerate}[label=(\alph*),itemsep=1pt]
          \item อัตราส่วนร่วม ($r$) = \ans{$\frac{6}{2} = 3$}
          \item พจน์ทั่วไป ($a_n$) = \ans{$2 \cdot 3^{n-1}$}
          \item พจน์ที่ 8 ($a_8$) = $2 \cdot 3^7 = 2(2,187) =$ \ans{$4,374$}
        \end{enumerate}
\end{enumerate}
\vspace{0.2em}

%% ---------- กิจกรรม 2 ----------
\noindent\textbf{กิจกรรมที่ 2 : อนุกรมและผลรวมสะสม (Series \& Summations)} \hfill \textit{(8 นาที)}
\begin{enumerate}[label=\arabic*)]
  \item \textbf{สัญกรณ์ซิกมา ($\sum$):} จงเขียนกระจายพจน์และหาผลลัพธ์ของ $\sum_{i=1}^4 (3i + 2)$\\[0.15em]
        กระจายพจน์: \ans{$(3(1)+2) + (3(2)+2) + (3(3)+2) + (3(4)+2) = 5 + 8 + 11 + 14$} = \ans{38}
  \item \textbf{ชาเลนจ์ออมเงิน 100 วัน (เทคนิคการจับคู่ของเกาส์):}\\
        วันที่ 1 ออม 1 บาท จนถึงวันที่ 100 ออม 100 บาท ($S = 1 + 2 + 3 + \dots + 100$)\\
        จงแสดงวิธีคิดด้วยสูตร $S_n = \frac{n}{2}(a_1 + a_n)$:\\[0.15em]
        $S_{100} = \frac{100}{2}(\ans{1} + \ans{100}) = 50 \times \ans{101} =$ \ans{5,050 บาท}
  \item \textbf{แผนการออมเงิน 20 วัน:} ออมวันแรก 5 บาท หยอดเพิ่มวันละ 3 บาท ($5 + 8 + 11 + 14 + \dots$)\\
        จงหาผลรวมเงินออม 20 วันแรก ($S_{20}$) ด้วยสูตร $S_n = \frac{n}{2}[2a_1 + (n - 1)d]$:\\[0.15em]
        $S_{20} = \frac{20}{2}[2(5) + (20-1)(3)] = 10[10 + \ans{57}] = 10 \times \ans{67} =$ \ans{670 บาท}
\end{enumerate}
\vspace{0.2em}

%% ---------- ท้าทายท้ายหน้า 1 ----------
\noindent\textbf{ท้าทายเช็กความเข้าใจ (หน้า 1):} เปรียบเทียบ \textbf{แผน ก} (เริ่ม 1,000 บาท ฝากเพิ่มเดือนละ 1,000 บาท: เลขคณิต) vs \textbf{แผน ข} (เริ่ม 1,000 บาท โตทบต้นเท่าตัวทุกเดือน $r = 2$: เรขาคณิต) ในเดือนที่ 10 แผนใดได้เงินมากกว่า และต่างกันกี่บาท?\\[0.15em]
\ans{แผน ก: $a_{10} = 1,000 + 9(1,000) = 10,000$ บาท \quad vs \quad แผน ข: $a_{10} = 1,000 \times 2^9 = 1,000 \times 512 = 512,000$ บาท}\\
\ans{ตอบ: แผน ข ได้เงินมากกว่าแผน ก อยู่ $512,000 - 10,000 = \mathbf{502,000\text{ บาท}}$ (มากกว่าถึง 51.2 เท่า!)}
\vspace{0.1em}

\clearpage
%% ====================== หน้าที่ 2 ======================
\noindent
\begin{minipage}[c]{0.78\textwidth}
  {\Large\bfseries\color{navy} [เฉลย] ใบกิจกรรมในชั้นเรียน บทที่ 3 : ความสัมพันธ์และฟังก์ชัน (สัปดาห์ที่ 9)}\\[0.1em]
  {\normalsize รายวิชา 4091606 คณิตศาสตร์สำหรับคอมพิวเตอร์ --- หน้าที่ 2}\\[0.5em]
  {\bfseries\color{ans} เฉลยกิจกรรมที่ 3, 4, 5 (Master Answer Key)}
\end{minipage}\hfill
\begin{minipage}[c]{0.2\textwidth}
  \centering
  \fbox{\parbox{2.8cm}{\centering\bfseries\color{navy}เฉลยประจำสัปดาห์\\ Week 09 KEY}}
\end{minipage}
\par\vspace{0.3em}\hrule\vspace{0.4em}

%% ---------- กิจกรรม 3 ----------
\noindent\textbf{กิจกรรมที่ 3 : อนุกรมเรขาคณิตจำกัดและอนันต์ (Finite \& Infinite Geometric Series)} \hfill \textit{(8 นาที)}
\begin{enumerate}[label=\arabic*)]
  \item \textbf{ผลรวมอนุกรมเรขาคณิตจำกัด:} แผนการออมเงินแบบทวีคูณ $3 + 6 + 12 + 24 + \dots$ จงหาผลรวม 6 วันแรก ($S_6$)\\[0.15em]
        วิธีทำ: $a_1 = 3, \ r = $ \ans{2} , \quad สูตร $S_6 = \frac{a_1(r^6 - 1)}{r - 1} =$ \ans{$\frac{3(2^6 - 1)}{2 - 1} = \frac{3(64 - 1)}{1} = 3 \times 63 =$} \ans{189 บาท}
  \item \textbf{การพิจารณาภาวะลู่เข้าและการลู่ออก (Convergence vs Divergence):} จงระบุอัตราส่วนร่วม $r$ และทำเครื่องหมาย $\checkmark$
        \begin{enumerate}[label=(\alph*),itemsep=2pt]
          \item การแบ่งเค้กทีละครึ่ง: $\frac{1}{2} + \frac{1}{4} + \frac{1}{8} + \dots \implies r = $ \ans{0.5} \quad \ans{[ $\checkmark$ ] ลู่เข้า (Convergent)} \quad [ \ ] ลู่ออก (Divergent)
          \item การออมเงินทวีคูณไม่จำกัด: $1 + 2 + 4 + 8 + \dots \implies r = $ \ans{2} \quad [ \ ] ลู่เข้า (Convergent) \quad \ans{[ $\checkmark$ ] ลู่ออก (Divergent)}
          \item อนุกรมแกว่งกวัด: $5 - 5 + 5 - 5 + \dots \implies r = $ \ans{-1} \quad [ \ ] ลู่เข้า (Convergent) \quad \ans{[ $\checkmark$ ] ลู่ออก (Divergent)}
        \end{enumerate}
  \item \textbf{ผลรวมอนุกรมเรขาคณิตอนันต์ ($S_\infty = \frac{a_1}{1 - r}$ เมื่อ $|r| < 1$):}
        \begin{enumerate}[label=(\alph*),itemsep=2pt]
          \item จงหาผลรวมอนันต์ของ $12 + 6 + 3 + 1.5 + \dots$ ($a_1 = 12, r = 0.5$)\\[0.1em]
                วิธีทำ: $S_\infty = \frac{a_1}{1 - r} =$ \ans{$\frac{12}{1 - 0.5} = \frac{12}{0.5} =$} \ans{24}
          \item จงแปลงทศนิยมซ้ำ $0.333\dots = \frac{3}{10} + \frac{3}{100} + \dots$ เป็นเศษส่วนอย่างต่ำ ($a_1 = 0.3, r = 0.1$)\\[0.1em]
                วิธีทำ: $S_\infty = \frac{0.3}{1 - 0.1} =$ \ans{$\frac{0.3}{0.9} = \frac{3}{9} =$} \ans{$\frac{1}{3}$}
        \end{enumerate}
\end{enumerate}
\vspace{0.2em}

%% ---------- กิจกรรม 4 ----------
\noindent\textbf{กิจกรรมที่ 4 : ฟังก์ชันเวียนเกิดในชีวิตประจำวัน (Recursive Functions in Daily Life)} \hfill \textit{(8 นาที)}
\begin{enumerate}[label=\arabic*)]
  \item \textbf{โครงสร้าง 2 ส่วนสำคัญ:} ฟังก์ชันเวียนเกิดประกอบด้วย 1) \ans{กรณีฐาน (Base Case)} และ 2) \ans{ขั้นตอนเวียนเกิด (Recursive Step)}\\
        หากไม่มีกรณีฐาน (Base Case) จะเกิดผลเสียอย่างไร: \ans{ฟังก์ชันจะเรียกตัวเองไม่รู้จบ ไม่สามารถหยุดการทำงานได้และหาคำตอบไม่ได้}
  \item \textbf{แบบจำลองดอกเบี้ยเงินฝากทบต้นรายปี ($10\%$ ต่อปี เงินต้น 10,000 บาท):}\\
        กรณีฐาน: $B(0) =$ \ans{10,000} \qquad ขั้นตอนเวียนเกิด: $B(n) =$ \ans{$1.10 \times B(n - 1)$}\\
        จงคำนวณ: $B(1) =$ \ans{11,000} บาท, \quad $B(2) =$ \ans{12,100} บาท, \quad $B(3) =$ \ans{13,310} บาท
  \item \textbf{แฟกทอเรียล ($n!$) กับการจัดลำดับสิ่งของ:} นิยาม $n! = n \times (n-1)!$ โดยมี Base Case $1! = 1$\\
        จงคลี่ขั้นตอนคำนวณ $4! = 4 \times 3! = 4 \times (3 \times 2!) = $ \ans{$4 \times (3 \times (2 \times 1)) =$} \ans{24}\\
        (หากมีคน 4 คน เข้าแถวซื้อของ จะจัดเรียงลำดับคิวได้ทั้งหมด \ans{24} แบบ)
  \item \textbf{ลำดับฟีโบนัชชีในธรรมชาติ ($F_n = F_{n-1} + F_{n-2}$ โดย $F_0 = 0, F_1 = 1$):}\\
        จงหาค่า: $F_2 = 1, \ F_3 = 2, \ F_4 = $ \ans{3}, \ $F_5 = $ \ans{5}, \ $F_6 = $ \ans{8}\\
        ยกตัวอย่างที่พบสัดส่วนฟีโบนัชชีในธรรมชาติ 1 ตัวอย่าง: \ans{กลีบดอกไม้ (ลิลลี่ 3, บัตเตอร์คัพ 5, เดซี่ 13, 21 กลีบ), ตาสับปะรด, ดอกทานตะวัน}
\end{enumerate}
\vspace{0.2em}

%% ---------- กิจกรรม 5 ----------
\noindent\textbf{กิจกรรมที่ 5 : การสรุปและสะท้อนคิด (Summary Reflection)} \hfill \textit{(4 นาที)}\\
จงอธิบายความแตกต่างสั้นๆ ระหว่าง \textbf{ลำดับ (Sequence)} และ \textbf{อนุกรม (Series)} ในมุมมองของการบันทึกสมุดบัญชีธนาคาร:\\[0.15em]
\ans{1. ลำดับ (Sequence): คือยอดเงินที่ฝากเพิ่มเข้ามาในแต่ละงวดเดี่ยวๆ เช่น งวดที่ 1 ฝาก 1,000, งวดที่ 2 ฝาก 1,500 ($a_n$)}\\
\ans{2. อนุกรม (Series): คือยอดเงินรวมสะสมทั้งหมดที่มีอยู่ในสมุดบัญชีตั้งแต่เดือนแรกจนถึงปัจจุบัน ($S_n = a_1 + a_2 + \dots + a_n$)}
\end{document}
